\chapter{Financial Foundations and Classical Models} \label{ch:finance}

\section{Financial Derivatives}\label{sec:derivatives}

In this section, we introduce the main ideas behind Mathematical Finance and the role played by Mathematics in this field. We first informally introduce the concept of \textit{derivative} and give some examples of the main quoted derivatives. The rigorous mathematical treatment of hedging and market completeness is addressed afterwards.

\subsection{The Concept of a Derivative}

In the context of modern mathematical finance, a financial derivative is formally defined as a \textit{contingent claim}. Its value is derived from---or ``contingent'' upon---the state of some other underlying variables.

Let us fix a filtered probability space $(\Omega, \mathcal{F}, (\mathcal{F}_t)_{t \in [0, T]}, \PP)$, where $\PP$ represents the physical (or real-world) probability measure, and the filtration $(\mathcal{F}_t)_{t \geq 0}$ models the flow of information over time up to a finite time horizon $T$. Mathematically, a financial derivative with maturity $T$ is an $\mathcal{F}_T$-measurable random variable, denoted by $\Payoff_T$. A few contracts, such as perpetual American options, carry an infinite time horizon ($T=\infty$); these fall outside our finite-horizon framework and form no core part of the derivatives market, so we exclude them.

This measurability condition implies that the value of the derivative at maturity is fully determined by the information available at time $T$. The function determining this value is often referred to as the \textit{payoff function}, denoted by $\PayoffFn$, so that $\Payoff_T = \PayoffFn(\Spot_T)$ for a claim written on the terminal spot. The key idea behind the subsequent theory is being able to assign \textit{fair prices} to those contingent claims, deduced from quoted (and liquid) market data.

The prices of such claims are not independent of one another. A small set of standardised, highly
liquid contracts carries enough information to determine the price of a great many bespoke ones, and
the principle enforcing this consistency is the absence of \textbf{arbitrage}: informally, the
impossibility of a strictly positive risk-free profit from zero net initial investment.



% For now, it will be convenient to ask ourselves the following question. For a fixed future date $T$, how much should earning 1 unit of currency (say 1 USD) be priced today? Likely, something below 1 present USD, and the quantity is likely to decay as $T$ increases. We will call this transformation the \textbf{discount factor} from $t$ to $T$, denoted by $D(t,T)$. It will be convenient to model the risk free rate $r(t)$ so that 
% \begin{equation}
% D(t,T) = e^{-\int_t^T r(s) ds}; \; \; r(t) = -\frac{\partial}{\partial T} \log D(t,T).
% \end{equation}

 

% We can think this option as leaving the payment in a deposit that pays the rates dictated by the central bank on each currency, typically FED for USD and ECB for EUR. In a hypothetical market with a constant risk-free interest rate $r$,  $D(t,T) = e^{-r(T-t)}$.

% A single example fixes the idea. Let a non-dividend-paying stock trade at $S_0$ and let a forward
% contract on it have delivery price $K > S_0e^{rT}$, with $r$ the continuous risk-free rate. Borrowing
% $S_0$, buying the stock and shorting the forward produces, at maturity, a delivery of the stock
% against $K$ and a repayment of $S_0e^{rT}$, for a certain profit of $K - S_0e^{rT} > 0$. Consistency
% therefore forces $K = S_0e^{rT}$.

One key observation is that pricing these contracts relies on the comparison of strategies
(\textit{hedging}), rather than on the expected value of the payoff under the physical measure
$\PP$. A single example fixes the idea. Let a non-dividend-paying stock trade at $\Spot_0$ and let a
forward contract on it have delivery price $\Strike$, with $\rf$ the continuous risk-free rate.
Borrowing $\Spot_0$, buying the stock and shorting the forward delivers, at maturity, the stock
against $\Strike$ and a repayment of $\Spot_0e^{\rf T}$, for a certain profit of
$\Strike - \Spot_0e^{\rf T}$; the mirror strategy locks in the opposite sign. Consistency therefore
forces $\Strike = \Spot_0e^{\rf T}$, and it does so without any model for how $\Spot$ moves.

Real markets impose frictions --- transaction costs, bid-ask spreads, asymmetric borrowing and
lending rates, taxes, and the credit risk borne by whoever must be paid at maturity, which grows
with the horizon --- and these widen the identity just derived into a band rather than a point. But
for the fungible assets and electronic markets we shall be concerned with, the band is tight, and it
is precisely the narrowness of that band which makes the mathematics worth doing. 

\subsubsection*{The Underlying Asset}
The process driving the derivative's value is the \textit{underlying asset}, an
$(\mathcal{F}_t)$-adapted process $S = (S_t)_{0\leq t\leq T}$; a derivative may depend on several. In
practice $S_t$ is the price of an equity or index, a bond or interest rate, a commodity such as oil, gold or electricity, or an exchange rate.
Derivatives are then categorised by the functional form of their terminal payoff $\Payoff_T$.

\paragraph{1. Forwards and Futures}
These instruments possess a payoff that is a linear function of the underlying asset price at maturity. A prime example is the \textbf{forward contract}, an agreement to buy or sell the asset at time $T$ for a pre-specified delivery price $K$. The payoff at time $T$ for the holder of a long position is:
\begin{equation}
    \Payoff_T = \Spot_T - \Strike
\end{equation}
Note that this payoff takes values in $[-\Strike,\infty)$ --- it is bounded below only because $\Spot_T \geq 0$ --- so the holder carries a genuine liability whenever the asset price falls below $\Strike$.

\paragraph{2. Options}
Options provide the \textit{right}, but not the \textit{obligation}, to transact. This asymmetry introduces convexity (non-linearity) into the payoff structure, truncating downside risk.
\begin{itemize}
    \item \textbf{European call:} the right to buy $\Spot_T$ at the strike $\Strike$. The payoff is defined as:
    \begin{equation}
        \Payoff_T = (\Spot_T - \Strike)^+ = \max(\Spot_T - \Strike, 0)
    \end{equation}
    \item \textbf{European put:} the right to sell $\Spot_T$ at the strike $\Strike$. The payoff is defined as:
    \begin{equation}
        \Payoff_T = (\Strike - \Spot_T)^+ = \max(\Strike - \Spot_T, 0)
    \end{equation}
        \item \textbf{American options:} American call or put options are similar to their European counterparts, with the key difference being that the option can be \textit{exercised} at any time up to maturity $T$.
\end{itemize}

Derivatives enable \textbf{hedging}, that is,  allow agents to transfer specific risks (e.g., volatility or tail risk) to those willing to bear them. In an idealized complete market, derivatives allow for the replication of any payoff profile. In a sense, liquid derivatives markets provide information to the market. Some market participants may be willing to join for pure speculation, while others genuinely need to hedge their positions, looking for a safer and more efficient way to transfer risk. Collectively, markets (and especially derivatives markets) act as an information network, where the price of a derivative reflects the market's consensus on the underlying asset's future distribution, given the payoffs involved.

These concepts will be rigorously defined in the next subsection. However, the main question remains: what price should these contracts hold? Before entering the mathematical machinery, we can establish bounds based on physical intuition, much like conservation laws:
\begin{itemize}
    \item If $T=0$, meaning we are pricing a claim that expires immediately (and only depends on the underlying at maturity, not on its trajectory), the price must equal the payoff. For a European call with strike $\Strike$ and underlying $\Spot_0$, its price must be $(\Spot_0-\Strike)^+$.
    \item If an asset gives its buyer strictly more optionality compared to another derivative, its price must be greater than or equal to the latter. For example, the price of an American call must be bounded below by the price of a European call with the same strike and maturity.
\end{itemize}

\subsection{Derivative Pricing}\label{subsec:derivativepricing}

We now turn to the rigorous problem of pricing contingent claims. While elementary treatments often assume asset prices follow geometric Brownian motions and interest rates are constant, these assumptions are insufficient for a general mathematical theory.

We work in a filtered probability space $(\Omega, \mathcal{F}, (\mathcal{F}_t)_{0 \le t \le T}, \PP)$ satisfying the usual conditions (right-continuity and completeness). The time horizon $T$ is finite.

Let the market consist of $n+1$ assets denoted by the vector process $S = (S^0, S^1, \dots, S^n)$. We assume that each component $S^i$ is an \textbf{Itô process}. While the most abstract fundamental theorems of asset pricing rely on the broader class of general semimartingales, Itô processes provide the precise mathematical framework required for the diffusion models prevalent in this thesis, seamlessly connecting with the stochastic calculus developed in Section \ref{sec:stochasticcalculus}.

We distinguish the $0$-th asset, $S^0$, as the \textbf{numeraire}.
\begin{definition}[Numeraire]
A numeraire is a price process $S^0$ that is strictly positive ($S^0_t > 0$ almost surely for all $t$) and non-dividend paying. It serves as the unit of account for the market.
\end{definition}
\noindent While $\Num$ is often chosen to be the bank account $\Num_t = \exp\bigl(\int_0^t \rf_s\,ds\bigr)$, the theory holds for \textit{any} valid Itô process numeraire (e.g., a foreign currency or a stock). The bank account is merely a convenient particular case because it explicitly models the time value of money.

This is the point at which to record two objects that will be used constantly and are easy to
conflate. Capitalising one unit of currency forward from $t$ to $T$ at the short rate multiplies it
by $\exp\bigl(\int_t^T \rf_s\,ds\bigr)$, which is $\Num_T/\Num_t$; \emph{discounting} is the inverse
operation, and we write
\begin{equation}\label{eq:discountfactor}
\Disc{t}{T} := \exp\Bigl(-\int_t^T \rf_s\,ds\Bigr) = \frac{\Num_t}{\Num_T}
\end{equation}
for the resulting \textbf{discount factor}. The \textbf{zero-coupon bond} $\Bond{t}{T}$ is a
different kind of object: it is the price, agreed and paid at $t$, of a contract delivering one unit
of currency at $T$. One is a calculation, the other a traded asset, and the two must fit together.
They do, and the pricing rule \eqref{eq:riskneutralpricing} below says exactly how: taking $\Num$ as
numeraire and the payoff $\Payoff_T \equiv 1$,
\begin{equation}\label{eq:bonddiscount}
\Bond{t}{T} = \Num_t\,\Ep{\QQ}\Bigl[\frac{1}{\Num_T}\,\Big|\,\Ft\Bigr] = \Ep{\QQ}\bigl[\Disc{t}{T}\mid\Ft\bigr].
\end{equation}
The bond is therefore the \emph{expected} discount factor, not the exponential of any single rate.
When $\rf$ is deterministic the expectation does nothing and the two coincide, $\Bond{t}{T} = \Disc{t}{T}$;
that is the case assumed throughout Sections \ref{sec:blackscholes} and \ref{sec:othermodels}, and it
is why a constant $\rf$ lets us write $e^{-\rf(T-t)}$ everywhere without further comment. When $\rf$ is
stochastic the two genuinely differ, discounting no longer factorises out of an expectation, and one
switches to the $T$-forward measure of Definition \ref{def:forwardprice} --- the measure in which the
bond, rather than the bank account, is the unit of account, so that discounting becomes a
multiplication by $\Bond{t}{T}$ again. 

We define the \textbf{discounted price process} $X$ as the prices of the primary assets expressed in units of the numeraire:
\begin{equation}
    X_t = \left( 1, \frac{S^1_t}{S^0_t}, \dots, \frac{S^n_t}{S^0_t} \right)
\end{equation}
Because the denominator $S^0_t$ is strictly positive almost surely, the transformation $f(x, y) = x/y$ is smooth. By a direct application of the multidimensional Itô's formula, the discounted price processes $X_t^i$ are themselves Itô processes.

A trading strategy is a predictable process $\pi = (\pi^0, \pi^1, \dots, \pi^n)$, where $\pi^i_t$ represents the number of units of asset $i$ held at time $t$. The letter $\Payoff$ is reserved for payoffs, so strategies are written $\pi$ throughout. The value of the portfolio at time $t$ is the scalar process:
\begin{equation}
    V_t(\pi) = \sum_{i=0}^n \pi^i_t S^i_t
\end{equation}

For the strategy to be economically meaningful, it must be \textbf{self-financing}. In the language of stochastic integration, this implies that changes in wealth are generated solely by the fluctuations of the asset prices, with no external injection or extraction of capital:
\begin{equation}
    dV_t(\pi) = \sum_{i=0}^n \pi^i_t \, dS^i_t \label{eq:selffinancing}
\end{equation}
When working with discounted prices $X$, the self-financing condition simplifies. Applying Itô's product rule, it can be shown that a strategy is self-financing if and only if the discounted wealth process $V_t/S^0_t$ satisfies:
\begin{equation}
    \frac{V_t(\pi)}{S^0_t} = \frac{V_0(\pi)}{S^0_0} + \int_0^t \sum_{i=1}^n \pi^i_u \, dX^i_u = \frac{V_0}{S^0_0} + (\pi \cdot X)_t
\end{equation}
where $(\pi \cdot X)$ denotes the vector stochastic integral.

In continuous time the naive definition of arbitrage is insufficient, because ``doubling strategies'' --- strategies whose intermediate losses are unbounded below --- can manufacture a certain profit from zero capital within a finite horizon. To preclude this we restrict the admissible strategies.
\begin{definition}[Admissible Strategy]\label{def:admissible}
A self-financing strategy $\pi$ is called $a$-admissible if its discounted value process is uniformly bounded from below by a constant $-a$ almost surely (i.e., $(\pi \cdot X)_t \ge -a$ for all $t$). A strategy is admissible if it is $a$-admissible for some $a \ge 0$.
\end{definition}

We can now say precisely what is to be excluded. Let $\mathcal{K}$ be the set of random variables representing the terminal wealth of all admissible strategies starting with zero capital:
\begin{equation}
    \mathcal{K} = \{ (\pi \cdot X)_T \mid \pi \text{ is admissible and } V_0(\pi) = 0 \}
\end{equation}
The market admits \textbf{no arbitrage} (NA) if $\mathcal{K}\cap L^{0}_{+} = \{0\}$; that is, no admissible strategy starting from nothing finishes non-negative and strictly positive with positive probability. This is the classical condition, and in continuous time it is too weak --- not because it tolerates an outright arbitrage, but because the set it constrains is not closed, so that a \emph{limit} of near-arbitrages may escape it. Delbaen and Schachermayer \cite{delbaenschachermayer1994} showed that the condition which is genuinely equivalent to the existence of an equivalent martingale measure is \textit{No Free Lunch with Vanishing Risk}. Writing $\mathcal{C}$ for the cone of claims dominated by elements of $\mathcal{K}$, that is $\mathcal{C} = \{ g \in L^\infty(\mathcal{F}_T) \mid \exists f \in \mathcal{K},\ g \le f \}$, the market satisfies \textbf{NFLVR} if:
\begin{equation}
    \overline{\mathcal{C}} \cap L_+^\infty = \{0\}
\end{equation}
where $\overline{\mathcal{C}}$ is the closure of $\mathcal{C}$ in the $L^\infty$ norm. Intuitively, this condition prevents not just exact arbitrage, but also sequences of strategies that approximate an arbitrage opportunity with risk converging to zero.

\vspace{0.2cm}

We are now equipped to state the core result linking market topology to measure theory.

\begin{definition}[Equivalent Local Martingale Measure]
A probability measure $\QQ$ is an Equivalent Local Martingale Measure (ELMM) for the numeraire $S^0$ if:
\begin{enumerate}
    \item $\QQ \sim \PP$ (they share the same null sets).
    \item The discounted price process $X_t = S_t / S^0_t$ is a \textbf{local martingale} under $\QQ$.
\end{enumerate}
\end{definition}

\begin{theorem}[First Fundamental Theorem of Asset Pricing]\label{thm:firstFTAP}
Assume the process $X$ is locally bounded. The market satisfies the NFLVR condition if and only if there exists at least one Equivalent Local Martingale Measure $\QQ$. \hfill $\square$
\end{theorem}

The other reasonable question to ask ourselves is whether every contingent claim can be replicated by a self-financing strategy. Before stating the second theorem, we must rigorously define market completeness.

\begin{definition}[Market Completeness]
A market is \textbf{complete} if every contingent claim $\Payoff_T \in L^\infty(\mathcal{F}_T)$ is \textit{attainable}. That is, for every bounded random variable $\Payoff_T$, there exists an admissible self-financing strategy $\pi$ and an initial capital $x \in \R$ such that:
\begin{equation}
    \frac{\Payoff_T}{S^0_T} = x + (\pi \cdot X)_T \quad \PP\text{-almost surely.}
\end{equation}
\end{definition}

Note that the statement is made under the physical measure $\PP$. This is deliberate: at this stage $\QQ$ has not been shown to be unique --- indeed its uniqueness is exactly what the next theorem characterises --- so quantifying over ``$\QQ$-almost surely'' would be ambiguous. Since every equivalent local martingale measure shares its null sets with $\PP$, the choice is harmless, but it is the honest one.

Two function spaces appear in this circle of ideas and it is worth saying once how they fit together. Completeness is defined above by quantifying over $L^\infty(\mathcal{F}_T)$, which is the Delbaen--Schachermayer convention and the right one for the Fundamental Theorems, boundedness being what makes their separating-hyperplane argument work. The Martingale Representation Theorem \ref{thm:mrt}, by contrast, is an $L^2$ statement, because it is at bottom an assertion about an orthogonal projection in a Hilbert space. The two are reconciled by the inclusion $L^\infty(\mathcal{F}_T) \subset L^2(\mathcal{F}_T)$, valid because $\PP$ is a finite measure: every bounded claim is square-integrable, so the representation applies to every claim the definition above quantifies over. We use whichever space the argument at hand requires, and this inclusion is the only thing needed to pass between them.

\begin{theorem}[Second Fundamental Theorem of Asset Pricing]\label{thm:secondFTAP}
Assume that the set of equivalent local martingale measures $\mathcal{M}^e$ is non-empty (i.e., the market is arbitrage-free). The market is complete if and only if $\mathcal{M}^e$ is a singleton set (i.e., $\QQ$ is unique). \hfill $\square$
\end{theorem}

If the market is arbitrage-free and complete (i.e., there exists a unique $\QQ$), the fair price at time $t$ of a generic derivative paying $\Payoff_T$ at maturity is given by the expectation of the discounted payoff:
\begin{equation}
    \frac{V_t}{S^0_t} = \Ep{\QQ} \left[ \frac{\Payoff_T}{S^0_T} \bigg| \mathcal{F}_t \right] \implies V_t = S^0_t \, \Ep{\QQ} \left[ \frac{\Payoff_T}{S^0_T} \bigg| \mathcal{F}_t \right] \label{eq:riskneutralpricing}
\end{equation}

This formulation reveals that ``risk-neutral probabilities'' are not absolute; they are  linked to the choice of numeraire.

Let $\widetilde{\Num}$ be a different numeraire (strictly positive semimartingale). Let $\QQ^{\Num}$ be the measure associated with numeraire $\Num$, and $\QQ^{\widetilde{\Num}}$ the measure associated with $\widetilde{\Num}$.
Prices must be consistent regardless of the unit of account. Therefore, for any asset $Z$:
\begin{equation}
    \frac{Z_t}{S^0_t} = \Ep{\QQ^{\Num}}\left[ \frac{Z_T}{S^0_T} \mid \mathcal{F}_t \right] \quad \text{and} \quad \frac{Z_t}{\widetilde{\Num}_t} = \Ep{\QQ^{\widetilde{\Num}}}\left[ \frac{Z_T}{\widetilde{\Num}_T} \mid \mathcal{F}_t \right]
\end{equation}

The change of measure from $\QQ^{\Num}$ to $\QQ^{\widetilde{\Num}}$ is defined by the Radon--Nikodym derivative, which turns out to be exactly the normalised ratio of the numeraires:
\begin{equation}
    \frac{d\QQ^{\widetilde{\Num}}}{d\QQ^{\Num}}\bigg|_{\mathcal{F}_T} = \frac{\widetilde{\Num}_T / S^0_T}{\widetilde{\Num}_0 / S^0_0} \label{eq:numerairechange}
\end{equation}
This elegantly shows that changing the numeraire is mathematically equivalent to performing a change of probability measure. This technique is particularly powerful in interest rate theory (e.g., switching to the $T$-forward measure to efficiently price bond options and swaptions).

\vspace{0.2cm}

Finally, we can now state the Feynman--Kac formula, which provides a powerful link between the PDE approach and the probabilistic representation of derivative prices.

\begin{theorem}[Feynman--Kac Formula]\label{thm:feynmankac}
    Let $X$ be an $n$-dimensional Itô diffusion process on the interval $[t, T]$ driven, under a measure $\QQ$, by the stochastic differential equation
    \begin{equation}
        dX_s = \mu(s, X_s)\,ds + \sigma(s, X_s)\,dW_s, \quad s \in [t, T],
    \end{equation}
    with initial condition $X_t = x$, where $W$ is a standard $n$-dimensional $\QQ$-Brownian motion, and the diffusion matrix is given by $\Sigma(s, x) = \sigma(s,x)\sigma(s,x)\transp$. Let $\PayoffFn: \R^n \rightarrow \R$ be a Borel-measurable terminal payoff function, and $\rf: [0,T] \times \R^n \rightarrow \R$ be a continuous function representing the continuously compounded interest rate. Assume $u(t,x)$ is a sufficiently regular function (typically $u \in C^{1,2}([0,T) \times \R^n) \cap C([0,T] \times \R^n)$) that satisfies the following linear parabolic partial differential equation:
    \begin{equation}\label{FK_PDE}
        \frac{\partial u}{\partial t}(t,x) + \sum_{i=1}^n \mu_i(t,x) \frac{\partial u}{\partial x_i}(t,x) + \frac{1}{2} \sum_{i,j=1}^n \Sigma_{ij}(t,x) \frac{\partial^2 u}{\partial x_i \partial x_j}(t,x) - \rf(t,x)u(t,x) = 0,
    \end{equation}
    for all $(t,x) \in [0,T) \times \R^n$, subject to the terminal condition:
    \begin{equation}
        u(T,x) = \PayoffFn(x).
    \end{equation}
    
    If $u(t,x)$ satisfies a suitable polynomial growth condition preventing the existence of explosive solutions, then $u(t,x)$ admits the stochastic representation:
    \begin{equation}\label{FK_Expectation}
        u(t,x) = \Ep{\QQ} \left[ \left. e^{-\int_t^T \rf(s, X_s) ds} \PayoffFn(X_T) \right| X_t = x \right].
    \end{equation}
\qedopen
\end{theorem}

The subscript on the expectation is not decoration. Equation \eqref{FK_Expectation} is an
expectation under \emph{the measure making $X$ solve the stated equation} --- that is, the measure
under which the drift is $\mu$ --- and nothing in the theorem says which measure that is. When the
theorem is applied to option pricing in Section \ref{sec:blackscholes}, the drift supplied is
$\rf\Spot$ rather than $\mu\Spot$, so the measure in question is the risk-neutral $\QQ$ and
emphatically not the physical $\PP$. For further details on the Feynman--Kac formula, see \cite{karatzasshreve1991}.

\section{The Black--Scholes Framework}\label{sec:blackscholes}

Having established the general theory of arbitrage-free pricing and martingale measures, we now specialise our market model to the most famous and widely used framework: the Black--Scholes--Merton model, usually called the Black--Scholes model, or simply the BS model.

\subsection{The Model}
\label{subsec:bs_model}

We assume a frictionless market: no transaction costs, no bid-ask spread, unrestricted short 
selling, and a constant risk-free rate $\rf$ at which one may borrow and lend freely. Under the 
\emph{physical} measure $\PP$ --- the measure describing what actually happens --- the underlying is 
modelled by a geometric Brownian motion,
\begin{equation} \label{eq:BS_SDE_P}
    d\Spot_t = \mu \Spot_t \, dt + \sigma \Spot_t \, d\Wt_t, \qquad \Spot_0 > 0,
\end{equation}
where $\mu \in \R$ is the expected rate of return, $\sigma>0$ the volatility, and $\Wt$ a 
$\PP$-Brownian motion. By Example \ref{example:blackscholessde} the solution is 
$\Spot_t = \Spot_0\exp\bigl((\mu-\tfrac12\sigma^2)t + \sigma \Wt_t\bigr)$, which is positive for all 
$t$ --- as a limited-liability equity should be.

It is worth saying plainly what $\sigma$ \emph{is}, since the word ``volatility'' will shortly be
doing several jobs at once. Inside the model it is the diffusion coefficient of the \emph{log}-price:
by Example \ref{example:blackscholessde}, $d\log \Spot_t = (\mu-\tfrac12\sigma^2)\,dt + \sigma\,d\Wt_t$,
so that $\Var[\log \Spot_t] = \sigma^2 t$ and $\sigma^2$ is the rate at which variance accumulates.
Outside the model it is a statistic: if prices $\Spot_0,\Spot_1,\dots,\Spot_n$ are observed at
spacing $\Delta$, the log-returns $R_i := \log(\Spot_i/\Spot_{i-1})$ are, under
\eqref{eq:BS_SDE_P}, independent and $\Normal\bigl((\mu-\tfrac12\sigma^2)\Delta,\ \sigma^2\Delta\bigr)$,
so $\sigma$ is the standard deviation of the log-returns \emph{per unit time} and is estimated by the
sample standard deviation of the $R_i$ divided by $\sqrt{\Delta}$. In the Black--Scholes world the
parameter and the statistic are the same number. Everything in Section \ref{sec:othermodels} follows
from the fact that, in the market, they are not.

It is essential to begin under $\PP$ and not under a risk-neutral measure. Doing so makes the 
central result of the theory visible rather than assumed, and that result is the following: 
\textbf{the price of the derivative does not depend on $\mu$.} Two investors who disagree completely 
about the expected return of the stock, but agree on its volatility, must agree on the price of 
every option written on it. This is not intuitive, it is not a modelling convention, and it is worth 
seeing emerge from the computation rather than being handed the risk-neutral dynamics with $\mu$ 
already replaced by $\rf$.

\subsubsection*{The hedging argument}

Consider a derivative whose value is a smooth function $V(t,\Spot)$ of time and the current spot. 
Form the portfolio
\[
\Pi_t = -V(t,\Spot_t) + \Delta_t \Spot_t,
\]
consisting of a short position in the derivative and $\Delta_t$ units of the underlying. Assuming 
the position is self-financing in the sense of \eqref{eq:selffinancing}, its value evolves as
\begin{equation}\label{eq:portfoliodynamics}
    d\Pi_t = -dV_t + \Delta_t \, d\Spot_t .
\end{equation}
Itô's rule (Theorem \ref{theorem:itorule}) applied to $V(t,\Spot_t)$, using 
$(d\Spot_t)^2 = \sigma^2\Spot_t^2\,dt$ from Table \ref{tab:ItoTableScalar}, gives
\begin{equation}\label{eq:dV}
    dV_t = \left( \frac{\partial V}{\partial t} + \mu \Spot_t\frac{\partial V}{\partial \Spot} 
    + \frac{1}{2} \sigma^2 \Spot_t^2 \frac{\partial^2 V}{\partial \Spot^2} \right) dt 
    + \sigma \Spot_t \frac{\partial V}{\partial \Spot}\, d\Wt_t .
\end{equation}
Substituting \eqref{eq:BS_SDE_P} and \eqref{eq:dV} into \eqref{eq:portfoliodynamics} and collecting 
terms,
\begin{equation}\label{eq:portfoliogrouped}
    d\Pi_t = \underbrace{\left( -\frac{\partial V}{\partial t} 
    - \frac{1}{2}\sigma^2 \Spot_t^2 \frac{\partial^2 V}{\partial \Spot^2} 
    + \mu \Spot_t\Bigl[\Delta_t - \frac{\partial V}{\partial \Spot}\Bigr] \right)}_{\text{drift}} dt 
    + \underbrace{\sigma \Spot_t\left( \Delta_t - \frac{\partial V}{\partial \Spot} \right)}_{\text{diffusion}} d\Wt_t .
\end{equation}
The portfolio is exposed to randomness only through the second term. Choosing
\begin{equation}\label{eq:delta}
    \Delta_t = \frac{\partial V}{\partial \Spot}(t,\Spot_t)
\end{equation}
annihilates it. This is the \emph{delta} of the derivative, and \eqref{eq:delta} is the hedge.

Now observe what \eqref{eq:portfoliogrouped} has become. The same choice \eqref{eq:delta} that kills 
the diffusion term also kills the bracket multiplying $\mu$ in the drift. \textbf{The expected return 
of the stock has cancelled}, and it has done so for a structural reason rather than by accident: 
$\mu$ enters $d\Spot_t$ and $dV_t$ through the identical channel, so a position that is neutral to 
the random part of $\Spot$ is automatically neutral to its drift as well. What remains,
\[
d\Pi_t = \left( -\frac{\partial V}{\partial t} - \frac{1}{2}\sigma^2 \Spot_t^2\frac{\partial^2 V}{\partial \Spot^2}\right) dt,
\]
is riskless. A riskless self-financing portfolio must earn the risk-free rate --- otherwise, by 
borrowing or lending against it, one manufactures an arbitrage in the sense of Subsection 
\ref{subsec:derivativepricing} --- so $d\Pi_t = \rf \Pi_t\,dt$ with 
$\Pi_t = -V + \Spot_t \partial V/\partial \Spot$. Equating the two expressions and rearranging yields 
the Black--Scholes equation.

\begin{theorem}[Black--Scholes--Merton]\label{thm:blackscholespde}
Under the dynamics \eqref{eq:BS_SDE_P}, the value $V(t,\Spot)$ of a European derivative with payoff 
$V(T,\Spot) = \PayoffFn(\Spot)$ satisfies the linear parabolic partial differential equation
\begin{equation} \label{eq:BS_PDE}
    \frac{\partial V}{\partial t} + \rf \Spot \frac{\partial V}{\partial \Spot} 
    + \frac{1}{2} \sigma^2 \Spot^2 \frac{\partial^2 V}{\partial \Spot^2} - \rf V = 0, 
    \qquad (t,\Spot)\in[0,T)\times(0,\infty),
\end{equation}
together with the terminal condition $V(T,\Spot)=\PayoffFn(\Spot)$.
\end{theorem}

\textbf{\textit{Proof:}} The hedging argument above. \hfill $\blacksquare$

\medskip

Equation \eqref{eq:BS_PDE} contains $\rf$ but not $\mu$, which is the claim made above, now proved.
It is also exactly the Feynman--Kac equation of Theorem \ref{thm:feynmankac} for a diffusion with 
drift $\rf \Spot$ and diffusion coefficient $\sigma \Spot$, whence the probabilistic representation
\begin{equation}\label{eq:bsexpectation}
V(t,\Spot) = e^{-\rf(T-t)}\,\Ep{\QQ}\bigl[\PayoffFn(\Spot_T)\,\big|\,\Spot_t = \Spot\bigr],
\qquad d\Spot_t = \rf \Spot_t\,dt + \sigma \Spot_t\,d\Wt^{\QQ}_t .
\end{equation}
The passage from \eqref{eq:BS_SDE_P} to the dynamics in \eqref{eq:bsexpectation} is precisely 
Girsanov's theorem with $\alpha = (\rf-\mu)/\sigma$, that is, with $-\alpha = (\mu-\rf)/\sigma$ the
\emph{market price of risk}: the measure
change subtracts a drift and leaves $\sigma$ untouched, exactly as discussed after Theorem 
\ref{girsanovtheorem}. That the hedging argument and the martingale argument give the same answer is 
not a coincidence but the content of Feynman--Kac, and the two viewpoints will alternate throughout 
what follows --- the PDE when we want to compute, the expectation when we want to reason.

Finally, the market is complete: by Theorem \ref{thm:mrt} and Remark \ref{rem:mrtpricing}, the 
filtration is generated by the single Brownian motion driving the single traded asset, so every 
square-integrable claim is attainable and $\QQ$ is unique. The delta \eqref{eq:delta} is the 
integrand $\varphi$ whose abstract existence the representation theorem guarantees; here we have it 
explicitly.

\subsubsection*{The closed-form solution}

For a European call, $\PayoffFn(\Spot)=(\Spot-\Strike)^+$, evaluating \eqref{eq:bsexpectation} 
against the log-normal density gives
\begin{equation}\label{eq:bscall}
    C_{\mathrm{BS}}(t,\Spot) = \Spot \, N(d_+) - \Strike e^{-\rf(T-t)} N(d_-), 
    \qquad d_{\pm} = \frac{\ln(\Spot/\Strike) + (\rf \pm \tfrac12\sigma^2)\tau}{\sigma\sqrt{\tau}},
\end{equation}
with $\tau = T-t$ and $N$ the standard normal distribution function. The corresponding put price 
follows from put--call parity, Proposition \ref{prop:putcallparity} below.

One might instead assume \textbf{normal} (Bachelier) dynamics $d\Spot_t = \mu\,dt + \sigma\,d\Wt_t$,
as in the original 1900 thesis \cite{bachelier1900}. This permits negative prices, which is why it was criticised for
equities, but it is standard for interest rates and for spread options where negative values are
economically meaningful. For equities the log-normal convention prevails, on account of the limited
liability of shareholders.

\subsubsection*{The Greeks}

The delta \eqref{eq:delta} was singled out above as the hedge, and identified in Remark
\ref{rem:mrtpricing} with the integrand $\varphi$ whose existence the Martingale Representation
Theorem guarantees. It is not the only sensitivity that matters, and since the whole of Subsection
\ref{subsec:smile_skew} is an argument about one of the others, we record them here. Write
$\tau = T-t$ and let $N'(x) = e^{-x^2/2}/\sqrt{2\pi}$ be the standard normal density.

\begin{proposition}[Black--Scholes sensitivities]\label{prop:greeks}
For the European call \eqref{eq:bscall} on a non-dividend-paying underlying,
\begin{align}
\Delta &:= \frac{\partial C}{\partial \Spot} = N(d_+), &
\Gamma &:= \frac{\partial^2 C}{\partial \Spot^2} = \frac{N'(d_+)}{\Spot\,\sigma\sqrt{\tau}}, \label{eq:deltagamma}\\[4pt]
\mathcal{V} &:= \frac{\partial C}{\partial \sigma} = \Spot\,N'(d_+)\sqrt{\tau}, &
\theta &:= \frac{\partial C}{\partial t} = -\frac{\Spot\,N'(d_+)\sigma}{2\sqrt{\tau}} - \rf \Strike e^{-\rf\tau}N(d_-). \label{eq:vegatheta}
\end{align}
\qedopen
\end{proposition}

The names are delta, gamma, vega and theta respectively. Two identities among these are worth more than the formulae themselves. The first is that
$\Gamma$ and $\mathcal{V}$ are proportional,
\begin{equation}\label{eq:vegagamma}
\mathcal{V} \;=\; \Spot^2\,\sigma\,\tau\,\Gamma ,
\end{equation}
as one verifies immediately from \eqref{eq:deltagamma} and \eqref{eq:vegatheta}. Convexity in the
spot and sensitivity to volatility are therefore the \emph{same} exposure, measured in two units ---
which is the analytic content of the informal statement that an option buyer is long volatility
precisely because the payoff is convex. The second is that substituting \eqref{eq:deltagamma} into the Black--Scholes equation
\eqref{eq:BS_PDE} rearranges it into
\begin{equation}\label{eq:thetagamma}
\theta + \tfrac12\sigma^2\Spot^2\,\Gamma \;=\; \rf\bigl(C - \Spot\,\Delta\bigr),
\end{equation}
which says that for a delta-hedged position ($\Delta$ neutralised, so that the right-hand side is
just the financing of the residual cash) the time decay $\theta$ and the convexity
$\Gamma$ are in exact opposition. This is the whole economics of running an option book in one line:
the gamma that makes the position profit from movement is paid for, day by day, by theta, and
\eqref{eq:thetagamma} fixes the exchange rate between them. Equation \eqref{eq:BS_PDE} is usually
read as a pricing equation; \eqref{eq:thetagamma} is the same equation read as a statement about a
hedged portfolio, and it is the form a trader would recognise.

Vega is the one to keep in view. It is strictly positive, so a call is worth more when volatility is
higher; and it is \emph{not} a sensitivity to a quantity the model admits to be uncertain --- in
\eqref{eq:BS_SDE_P} the parameter $\sigma$ is a constant. Differentiating a price with respect to a
constant is a symptom that the model is being used outside its own assumptions, and the entire
subject of Section \ref{sec:othermodels} is the attempt to build models in which the quantity vega
measures is a genuine state variable rather than a parameter one perturbs by hand. Subsection
\ref{subsec:smile_skew} is where that symptom becomes impossible to ignore.

\subsection{Dividends, and Why the Forward is the Right Variable}\label{subsec:dividends}

The derivation above quietly assumed that holding the stock produces no cash flows. Real equities pay 
dividends, and the way they do so is awkward enough to force a change of variable that we shall keep 
for the rest of the dissertation.

Suppose first that dividends are paid \emph{continuously} at a constant proportional rate $\divy$. 
The holder of the stock receives $\divy \Spot_t\,dt$ over $[t,t+dt]$, so the self-financing condition 
\eqref{eq:portfoliodynamics} acquires the extra term $\Delta_t \divy \Spot_t\,dt$; repeating the 
argument verbatim replaces the coefficient $\rf$ of $\partial V/\partial \Spot$ in \eqref{eq:BS_PDE} 
by $\rf - \divy$, and $\rf$ by $\rf-\divy$ in $d_\pm$. Nothing structural changes --- but the closed 
form does not follow from those two substitutions alone, and it is worth writing out, since applying 
them mechanically to \eqref{eq:bscall} gives the wrong price. The discounted-spot term acquires its 
own factor:
\begin{equation}\label{eq:bscalldiv}
C(\Spot,\tau) = \Spot e^{-\divy\tau}N(d_+) - \Strike e^{-\rf\tau}N(d_-), \qquad
d_\pm = \frac{\log(\Spot/\Strike) + (\rf-\divy \pm \tfrac12\sigma^2)\tau}{\sigma\sqrt{\tau}} .
\end{equation}
The reason is visible in the Black-76 form \eqref{eq:black76}: what enters the formula is the forward 
$\Fwd = \Spot e^{(\rf-\divy)\tau}$, and substituting it into $e^{-\rf\tau}[\Fwd N(d_+) - \Strike N(d_-)]$ 
produces $e^{-\divy\tau}$ on the spot term automatically. This is a further argument for treating 
\eqref{eq:black76} rather than \eqref{eq:bscall} as the primitive object: the dividend correction is 
invisible there because the forward has already absorbed it. Continuous proportional dividends are,
however, the case usually presented and almost never the case observed: real equities pay discrete
cash amounts, and it is to those that we now turn.

Actual equity dividends are \emph{discrete cash amounts}, announced in advance and paid on known 
dates. If an amount $D_i$ is paid at time $t_i$, then no-arbitrage forces the spot to drop by that 
amount across the payment:
\begin{equation}\label{eq:dividendjump}
\Spot_{t_i} = \Spot_{t_i^-} - D_i .
\end{equation}
This is genuinely destructive of the framework above, and it is worth being precise about why. 
Equation \eqref{eq:dividendjump} says $\Spot$ has a \emph{predictable jump}: the process is not 
continuous, so Itô's rule as stated does not apply across $t_i$. Worse, the jump is \emph{additive} 
while geometric Brownian motion is multiplicative, so the two are structurally incompatible --- one 
cannot write $\Spot_t = \Spot_0 \mathcal{E}(\cdot)$ for a positive $\mathcal{E}$ and reproduce 
\eqref{eq:dividendjump}. Worse still, a fixed cash dividend can in principle exceed the spot price, 
so a model treating $D_i$ as deterministic does not even guarantee positivity. And since 
$\sigma$ multiplies $\Spot_t$ in \eqref{eq:BS_SDE_P}, a discontinuous drop in $\Spot$ produces a 
discontinuous drop in the absolute volatility, which is not what is observed.

The resolution is to stop modelling the spot and model the \emph{forward} instead.

\begin{definition}[Forward price]\label{def:forwardprice}
The \textbf{forward price} $\Fwd(t,T)$ is the delivery price agreed at $t$ for a transaction in the 
underlying at $T$ which has zero value at inception. Equivalently, writing $\Bond{t}{T}$ for the 
price at $t$ of a zero-coupon bond maturing at $T$,
a forward contract struck at $K$ has time-$t$ value
\begin{equation}\label{eq:forwardvalue}
V_t(K) = \Bond{t}{T}\;\Ep{\QT}\bigl[\Spot_T - K \mid \Ft\bigr],
\end{equation}
where $\QT$ is the $T$-forward measure associated with the numeraire $\Bond{\cdot}{T}$. Since 
$\Bond{t}{T} > 0$, the requirement that the contract be costless at inception, $V_t(\Fwd(t,T)) = 0$, 
determines the delivery price uniquely:
\begin{equation}\label{eq:forwarddef}
\Fwd(t,T) = \Ep{\QT}\bigl[\Spot_T \mid \Ft\bigr].
\end{equation}
\end{definition}

Identity \eqref{eq:forwarddef} says that the forward price is a $\QT$-martingale --- it is a 
conditional expectation --- and this is the whole point. It holds by construction, with no assumption 
whatever on interest rates, on dividends, or on their nature. In particular we deliberately do 
\emph{not} write $\Fwd(t,T) = \Spot_t e^{(\rf-\divy)(T-t)}$, which is valid only for constant $\rf$, 
constant $\divy$, and purely proportional dividends. Instead we record the relation abstractly, as an
\textbf{affine} one:
\begin{equation}\label{eq:spotforwardfactor}
\Fwd(t,T) = a_t\,\Spot_t + b_t, \qquad a_t > 0.
\end{equation}
 Both coefficients are adapted and
model-independent: $a_t$ is the \emph{proportional} part, assembled from the discount curve, the repo
or borrow cost and any proportional dividend yield, while $b_t$ is the \emph{cash} part, carrying the
announced discrete dividends. Their precise construction is a market-data question, not a modelling
one. What matters here is that \eqref{eq:spotforwardfactor} remains
valid for time-dependent $\rf(t)$ and $\divy(t)$, for discrete cash dividends, for proportional
dividends, and for any mixture of them; the purely proportional case of \eqref{eq:bscalldiv} is
$b_t = 0$. Modelling the forward as a driftless log-normal martingale,
\begin{equation}\label{eq:blackdynamics}
d\Fwd_t = \sigma \Fwd_t \, d\Wt^{T}_t,
\end{equation}
therefore has three advantages simultaneously. The dividend discontinuities are absorbed into
$b_t$ and never touch the dynamics, which stay diffusive. The drift vanishes, because a forward 
is already a martingale under its own measure --- so there is no drift to specify and none to get 
wrong. And the resulting formula is calibrated directly against quantities the market quotes.

\subsubsection*{The forward measure and Black-76}
Changing numeraire from the bank account to the zero-coupon bond $\Bond{\cdot}{T}$, in the manner of 
\eqref{eq:numerairechange}, converts \eqref{eq:bsexpectation} into an expectation of the undiscounted 
payoff under $\QT$, multiplied by the bond price $\Bond{t}{T}$, which under the deterministic
rates assumed here coincides with the discount factor \eqref{eq:discountfactor}. With the dynamics 
\eqref{eq:blackdynamics} this gives Black's 1976 formula.

\begin{proposition}[Black-76 \cite{black1976}]\label{prop:black76}
Under \eqref{eq:blackdynamics}, the time-$t$ value of a European call with strike $\Strike$ and 
maturity $T$ is
\begin{equation}\label{eq:black76}
C(t) = \Bond{t}{T}\Bigl[\Fwd_t\,N(d_+) - \Strike\,N(d_-)\Bigr], \qquad 
d_\pm = \frac{\ln(\Fwd_t/\Strike) \pm \tfrac12 \sigma^2\tau}{\sigma\sqrt{\tau}}, \quad \tau = T-t.
\end{equation}
\qedopen
\end{proposition}

Formula \eqref{eq:black76} is formally identical to \eqref{eq:bscall} with $\Spot$ replaced by 
$\Fwd_t$ and the interest rate removed from $d_\pm$ --- but note that the bond price 
$\Bond{t}{T}$ remains, multiplying the whole bracket. It has not disappeared; it has merely been 
collected outside. 

\subsubsection*{Put--call parity}

Before turning to implied volatility we record the one relation in this subject that holds 
independently of any model, and which we shall invoke repeatedly.

\begin{proposition}[Put--call parity]\label{prop:putcallparity}
Let $C(t)$ and $P(t)$ be the time-$t$ prices of European call and put options on the same underlying, 
with the same strike $\Strike$ and maturity $T$. Then, in any arbitrage-free market,
\begin{equation}\label{eq:putcallparity}
C(t) - P(t) = \Bond{t}{T}\bigl(\Fwd(t,T) - \Strike\bigr).
\end{equation}
\end{proposition}

\textbf{\textit{Proof:}} The payoffs satisfy the pointwise identity 
$(\Spot_T-\Strike)^+ - (\Strike-\Spot_T)^+ = \Spot_T - \Strike$, valid for every value of $\Spot_T$. 
Taking $\Ep{\QT}[\,\cdot\mid\Ft]$ and multiplying by $\Bond{t}{T}$, the right-hand side becomes 
$\Bond{t}{T}(\Fwd(t,T)-\Strike)$ by \eqref{eq:forwarddef}. \hfill $\blacksquare$

The proof used no dynamics --- not \eqref{eq:blackdynamics}, not even continuity of paths --- only 
the linearity of expectation and an algebraic identity between payoffs. Consequently 
\eqref{eq:putcallparity} must hold in every model in this dissertation, and it is the standard 
sanity check on any implementation: a pricer that violates put--call parity is wrong, and no appeal 
to model choice can rescue it.

Two consequences are worth noting now. First, since $C-P$ is model-free, calls and puts of the same 
strike carry \emph{identical} information about volatility; a call and a put struck at $\Strike$ must 
therefore have the same implied volatility, and the smile is a single function of strike rather than 
two. Second, \eqref{eq:putcallparity} is how the forward is extracted from quoted option prices in 
practice: the call-minus-put spread is linear in $\Strike$ with slope $-\Bond{t}{T}$ and root 
$\Strike = \Fwd(t,T)$, so a regression across strikes recovers both the bond price and the 
forward without any assumption about dividends. This is exactly how the coefficients $a_t$ and $b_t$ of
\eqref{eq:spotforwardfactor} are obtained: from market data, rather than postulated.

\subsection{Implied Volatility and the Inversion Problem}\label{subsec:impliedvol}
In professional markets, participants do not ``price'' options using Black--Scholes to find the dollar value. The market price $C^{\mathrm{mkt}}$ is determined by supply and demand. Instead, Black--Scholes is used as a translation mechanism. Since vega is strictly positive (Proposition \ref{prop:greeks}), the pricing function $\sigma \mapsto C_{\mathrm{BS}}(\sigma)$ is strictly \emph{increasing}, and it sweeps out the whole no-arbitrage interval as $\sigma$ runs over $(0,\infty)$; any quote inside that interval therefore has one and only one preimage $\ivol$:
\begin{equation}
    C_{\mathrm{BS}}(\Spot, \Strike, T, \rf, \divy, \ivol) = C^{\mathrm{mkt}},
\end{equation}
the dividend yield $\divy$ entering through \eqref{eq:bscalldiv}.
This 
implied volatility is the true ``price'' quoted in the market. It is a quoting convention rather than a forecast: what it summarises is the market's view of the root-mean-square volatility over the life of the option \emph{under the pricing measure} $\QQ$, which is not the same as its expectation under $\PP$.

\subsection{Smile and Skewness}
\label{subsec:smile_skew}

If the assumptions of Black--Scholes were satisfied physically (i.e., if returns were truly log-normal with constant variance), the implied volatility $\ivol$ would be constant across all strike prices $K$ and maturities $T$.

Empirical data reveals this is not the case. When we plot $\ivol(K)$ against $K$, we observe distinct shapes:
\begin{itemize}
    \item \textbf{The Skew (Smirk):} Typical in equity markets. $\ivol$ decreases as $K$ increases. This reflects ``Crash Phobia''---investors bid up the price of low-strike Puts (insurance against market drops), implying higher volatility for downside strikes.
    \item \textbf{The Smile:} Typical in Forex markets. $\ivol$ is lowest at the money ($\Strike \approx \Fwd$) and rises for strikes far from the money in \emph{both} directions. This reflects ``fat tails'' (leptokurtosis) in the return distribution: extreme moves in either direction are more probable than the Gaussian model predicts.
\end{itemize}

The two shapes are not merely descriptive. Both say that the market's risk-neutral distribution 
is not log-normal, and the following classical result makes the connection exact.

\begin{proposition}[Breeden--Litzenberger \cite{breeden1978}]\label{prop:breedenlitzenberger}
Let $C(\Strike,T)$ denote the price of a European call of strike $\Strike$ and maturity $T$, and 
suppose the risk-neutral law of $\Spot_T$ admits a density $\psi_T$. Then
\begin{equation}\label{eq:breedenlitzenberger}
    \psi_T(\Strike) = \frac{1}{\Bond{0}{T}}\,\frac{\partial^2 C}{\partial \Strike^2}(\Strike,T).
\end{equation}
\end{proposition}

\textbf{\textit{Proof:}} Write $C(\Strike,T) = \Bond{0}{T}\int_{\Strike}^{\infty}(x-\Strike)\psi_T(x)\,dx$. 
Differentiating once under the integral sign, the boundary term vanishes because the integrand is 
zero at $x=\Strike$, leaving 
$\partial_\Strike C = -\Bond{0}{T}\int_\Strike^\infty \psi_T(x)\,dx$. Differentiating a second time 
gives \eqref{eq:breedenlitzenberger}. \hfill $\blacksquare$

The result is remarkable for what it does not assume: no model, no dynamics, only that prices are 
discounted expectations of the payoff. A continuum of option quotes at a fixed maturity therefore 
determines the entire marginal distribution of $\Spot_T$ --- the market does not merely express an 
opinion about volatility, it expresses a complete probability law, and the smile is the visible 
trace of that law failing to be log-normal.

Two consequences of \eqref{eq:breedenlitzenberger} will be used constantly. First, since 
$\psi_T \geq 0$, arbitrage-free call prices must be \emph{convex} in strike; a violation is a 
butterfly arbitrage, executable by buying two wing options and selling two at the middle strike. 
Second, the formula involves a \emph{second} derivative of market data, which is numerically 
delicate in a way that Subsection \ref{subsec:localvol} will have to confront directly.

% \figurehere{Empirical implied volatility smile for a liquid index at two maturities, showing the 
% flattening of the skew with increasing $T$. To be produced from the market data of Chapter 
% \ref{ch:applications}.}

\section{Other Important Models}\label{sec:othermodels}

To account for the volatility smile and skew, we must extend the Black--Scholes framework. There are two primary approaches: making volatility a deterministic function of the state (Local Volatility) or making it a stochastic process (Stochastic Volatility).

\subsection{Local Volatility Models}\label{subsec:localvol}

The Local Volatility (LV) framework, introduced by Dupire \cite{dupire1994}, posits that volatility is a deterministic function of the current asset price and time: $\lvol = \lvol(t,\Spot_t)$. The dynamics under the risk-neutral measure are:
\begin{equation}\label{eq:localvolspot}
    d\Spot_t = (\rf-\divy)\,\Spot_t\,dt + \lvol(t,\Spot_t)\,\Spot_t\,d\Wt_t
\end{equation}
This model preserves the completeness of the market (there is only one source of randomness $\Wt_t$) while allowing the model to fit \textit{any} observed arbitrage-free implied volatility surface exactly.

\begin{convention}[Arguments of $\lvol$]\label{conv:lvolarguments}
The same function $\lvol$ is read in two coordinate systems, and we fix the convention once. In
\emph{state} coordinates it is a diffusion coefficient and its arguments are written time first,
$\lvol(t,\Spot)$, exactly as for the $\mu(t,x)$ and $\sigma(t,x)$ of Section
\ref{sec:stochasticcalculus}. In \emph{surface} coordinates it is a function of a contract rather
than of a state, and we follow the market in writing the maturity last: $\lvol(\Strike,T)$ against
strike, or $\lvol(\lmon,T)$ against the log-moneyness $\lmon$ of \eqref{eq:totalvariance}. The two
readings are the same object --- Theorem \ref{thm:dupire} evaluates the state-coordinate function
along $\Spot = \Strike$, $t = T$ --- and the order of the arguments is what says which reading is
meant. The implied volatility $\ivol$ and the total implied variance $\tvar$ exist only in surface
coordinates and are always written in the second form.
\end{convention}

\subsubsection*{The Dupire formula and Gatheral's transformation}

Dupire's insight was to link the unknown function $\lvol(\Strike,T)$ directly to observed market 
option prices. Applying the Fokker--Planck equation \eqref{eq:fokkerplanck} to the transition density 
of the diffusion and combining it with Breeden--Litzenberger yields the following.

\begin{theorem}[Dupire]\label{thm:dupire}
Suppose call prices $C(\Strike,T)$ are known for a continuum of strikes and maturities and are 
consistent with a diffusion of the form $d\Spot_t = (\rf-\divy)\Spot_t\,dt + \lvol(t,\Spot_t)\Spot_t\,d\Wt_t$. 
Then the local volatility function is uniquely determined by
\begin{equation} \label{eq:dupire}
    \lvol^2(\Strike, T) = \frac{\dfrac{\partial C}{\partial T} + (\rf-\divy)\,\Strike \dfrac{\partial C}{\partial \Strike} + \divy\,C}{\tfrac{1}{2} \Strike^2 \dfrac{\partial^2 C}{\partial \Strike^2}} .
\end{equation}
\qedopen
\end{theorem}

The content of \eqref{eq:dupire} is striking and deserves to be stated plainly: \emph{there is 
exactly one} local volatility diffusion consistent with a given arbitrage-free surface. The 
calibration problem has a unique solution, not a family of them, and the model can therefore 
reproduce every European quote simultaneously and exactly. Note also that the denominator is 
proportional to the risk-neutral density by Proposition \ref{prop:breedenlitzenberger}, so 
\eqref{eq:dupire} is well posed precisely when that density is strictly positive.

\subsubsection*{Static arbitrage and the implied-variance parameterisation}

Formula \eqref{eq:dupire} is, however, close to unusable as written. It requires one derivative in 
maturity and two in strike of raw market data, and market data are quoted on a sparse grid, carry 
bid-ask spreads, and are contaminated by microstructure noise. Numerical differentiation amplifies 
all three, and the practical symptom is a negative denominator --- that is, a negative local 
variance, which is meaningless.

The remedy is not better differentiation but a change of variables, together with an interpolation 
that is guaranteed arbitrage-free by construction. Define the \textbf{log-moneyness} and the 
\textbf{total implied variance}
\begin{equation}\label{eq:totalvariance}
\lmon := \log\frac{\Strike}{\Fwd(t,T)}, \qquad \tvar(\lmon,T) := \ivol^2(\lmon,T)\,T .
\end{equation}
Total variance is the natural coordinate because it, rather than implied volatility itself, is what 
accumulates linearly in time under Black--Scholes. We now state the conditions under which a surface 
$\tvar$ is free of \emph{static} arbitrage --- that is, admits no arbitrage exploitable by a 
buy-and-hold portfolio of European options, with no trading in between.

\begin{proposition}[Absence of static arbitrage]\label{prop:staticarbitrage}
A total implied variance surface $\tvar(\lmon,T)$ is free of static arbitrage if the following hold.
\begin{itemize}
    \item[i)] (\textbf{Calendar spread}) For each fixed $\lmon$, the map $T \mapsto \tvar(\lmon,T)$ 
    is non-decreasing:
    \begin{equation}\label{eq:calendarspread}
    \frac{\partial \tvar}{\partial T}(\lmon,T) \geq 0 .
    \end{equation}
    \item[ii)] (\textbf{Butterfly}) For each fixed $T$, the function
    \begin{equation}\label{eq:butterfly}
    g(\lmon) := \left(1 - \frac{\lmon\,\partial_\lmon \tvar}{2\tvar}\right)^{\!2} 
    - \frac{(\partial_\lmon \tvar)^2}{4}\left(\frac{1}{\tvar} + \frac{1}{4}\right) 
    + \frac{\partial^2_{\lmon\lmon} \tvar}{2}
    \end{equation}
    satisfies $g(\lmon) \geq 0$ for all $\lmon$.
    \item[iii)] (\textbf{Wings}) $\tvar(\lmon,T)$ grows at most linearly in $|\lmon|$ as 
    $|\lmon|\to\infty$; more precisely, the asymptotic slopes are bounded by $2$, which is Lee's 
    moment formula \cite{lee2004}.
\end{itemize}
\hfill $\square$
\end{proposition}

Each condition has a direct financial meaning. Condition i) says a longer-dated option cannot be 
cheaper than a shorter-dated one of the same moneyness --- otherwise one buys the long and sells the 
short. Condition ii) is convexity of the call price in strike, expressed in the new coordinates, and 
its violation is a butterfly spread with negative cost and non-negative payoff. Condition iii) 
prevents the extrapolated wings from implying infinite moments of $\Spot_T$.

In practice one first fits a parametric form guaranteeing i)--iii) --- Gatheral's SVI 
parameterisation \cite{gatheral2014} being the standard choice, with SABR common in rates markets --- and only then 
differentiates. Since the fitted surface is smooth and arbitrage-free by construction, the 
differentiation is stable and the local variance is automatically non-negative. This yields 
Gatheral's form of the Dupire formula, which is the version actually implemented.

\begin{theorem}[Dupire's formula in implied-variance coordinates]\label{thm:gatheral}
With $\lmon$ and $\tvar$ as in \eqref{eq:totalvariance},
\begin{equation}\label{eq:gatheral}
\lvol^2(\lmon,T) 
= \frac{\dfrac{\partial \tvar}{\partial T}}
{\;1 - \dfrac{\lmon}{\tvar}\dfrac{\partial \tvar}{\partial \lmon} 
+ \dfrac{1}{4}\left(-\dfrac{1}{4} - \dfrac{1}{\tvar} + \dfrac{\lmon^2}{\tvar^2}\right)\!\left(\dfrac{\partial \tvar}{\partial \lmon}\right)^{\!2} 
+ \dfrac{1}{2}\dfrac{\partial^2 \tvar}{\partial \lmon^2}\;} .
\end{equation}
\qedopen
\end{theorem}

Formula \eqref{eq:gatheral} is worth a closer look, because the two no-arbitrage conditions of 
Proposition \ref{prop:staticarbitrage} are visible in it. The numerator is exactly the calendar 
spread quantity \eqref{eq:calendarspread}, and the denominator is precisely the butterfly quantity 
$g(\lmon)$ of \eqref{eq:butterfly}. Hence
\[
\frac{\partial \tvar}{\partial T} \geq 0 \ \text{ and } \ g(\lmon) \geq 0 
\;\implies\; \lvol^2 \geq 0 .
\]

% so that \textbf{an implied volatility surface free of static arbitrage yields a non-negative local 
% variance}. The implication runs one way only, and it is worth being precise about why: a quotient is 
% non-negative whenever numerator and denominator share a sign, so $\lvol^2 \geq 0$ is also compatible 
% with both being negative, a configuration which certainly does admit arbitrage. What is true is the 
% chain
% \[
% \text{no static arbitrage} \implies \frac{\partial \tvar}{\partial T} \geq 0 \ \wedge \ g(\lmon)\geq 0 
% \implies \lvol^2 \geq 0 ,
% \]
% in which each of the two conditions is \emph{separately} necessary for absence of arbitrage --- 
% $\partial_T \tvar \geq 0$ for calendar spreads, $g(\lmon)\geq 0$ for butterflies --- and only the 
% last arrow fails to reverse.

% That the denominator of \eqref{eq:gatheral} is $g(\lmon)$ exactly, and not merely something of the 
% same sign, is the substantive observation here, and it is what makes the two-step procedure above 
% more than a matter of taste: arbitrage-freeness is not an aesthetic requirement imposed on the 
% interpolation, it is the condition under which \eqref{eq:gatheral} is guaranteed to return a usable 
% answer. It is also the reason this parameterisation, rather than the raw form \eqref{eq:dupire}, is 
% the one calibrated in practice.

% \subsubsection*{The limitation}

Local volatility fits the surface perfectly, and by Theorem \ref{thm:dupire} it is the unique 
diffusion that does. It also preserves market completeness, since a single Brownian motion drives a 
single traded asset and Theorem \ref{thm:mrt} applies unchanged. These are substantial virtues.

The defect is equally fundamental, and it concerns \emph{dynamics} rather than fit. In a local 
volatility model the instantaneous volatility is a deterministic function of the spot, so the future 
smile is entirely determined by where the spot goes. The consequence, observed universally in 
practice, is that the model predicts the smile flattening as the spot moves, whereas real smiles 
translate with the spot and retain their shape. A model may therefore reproduce every European price 
today and still be wrong about every European price tomorrow. This matters not at all for pricing 
European options --- their value depends only on the terminal marginal, which is fitted exactly --- 
and a great deal for anything path-dependent: barriers, cliquets, forward-starting options and 
autocallables all depend on the joint law across dates, which local volatility gets wrong. That 
observation is what motivates the next subsection.

\subsection{Stochastic Volatility Models}\label{subsec:stochvol}

Throughout this subsection, and in Subsection \ref{subsec:american} which follows it, interest rates and dividend yields are taken to be \emph{deterministic}. The assumption is worth stating once rather than repeating: it is what makes the $T$-forward measure coincide with the risk-neutral one, $\QT = \QQ$, so that expectations below may be written under $\QQ$ while still exploiting the $\QT$-martingale property of the forward; and it is what allows the discounting term in the American pricing operator of \eqref{eq:LCP_American} to be a deterministic function of time. With stochastic rates both statements require the change of numeraire of Subsection \ref{subsec:dividends} to be carried through explicitly.

To capture realistic dynamics of the volatility surface itself, we introduce a second source of randomness. In Stochastic Volatility (SV) models, the variance $v_t = \sigma_t^2$ follows its own Stochastic Differential Equation. Under the risk-neutral measure $\QQ$, the system is described by:
\begin{equation}
    \begin{cases}
        d\Spot_t = (\rf-\divy)\,\Spot_t\,dt + \sqrt{v_t}\,\Spot_t\,d\Wt_t^{\Spot} \\
        dv_t = \mu_v(v_t, t)\,dt + \sigma_v(v_t, t)\,d\Wt_t^{v}.
    \end{cases}
\end{equation}
 A critical parameter is the correlation $\rho$ between the two Brownian motions: $d\langle \Wt^{\Spot}, \Wt^{v} \rangle_t = \rho\,dt$.
\begin{itemize}
    \item If $\rho < 0$, a drop in price ($\Spot \downarrow$) is associated with an increase in volatility ($v \uparrow$), naturally generating the \textbf{skew} observed in equity markets.
    \item The ``vol of vol'' parameter (in the diffusion coefficient $\sigma_v$) controls the convexity or \textbf{curvature} of the smile.
\end{itemize}

\subsubsection*{The Heston model}

The most prominent SV model is that of Heston \cite{heston1993}, in which the variance follows the 
Cox--Ingersoll--Ross square-root process already analysed in Example \ref{example:cirwellposed}:
\begin{equation}\label{eq:heston}
    \begin{cases}
        d\Spot_t = (\rf-\divy)\,\Spot_t\,dt + \sqrt{v_t}\,\Spot_t\,d\Wt_t^{\Spot},\\[2pt]
        dv_t = \kappa (\theta - v_t)\,dt + \xi \sqrt{v_t}\,d\Wt_t^{v},
    \end{cases}
    \qquad d\bigl\langle \Wt^{\Spot}, \Wt^{v}\bigr\rangle_t = \rho\,dt,
\end{equation}
where $\kappa$ is the speed of mean reversion, $\theta$ the long-run variance, $\xi$ the volatility 
of volatility, and $\rho$ the correlation, understood in the sense of Definition 
\ref{def:correlatedbm}. We recall from Subsection \ref{subsec:sdes} that \eqref{eq:heston} is well 
posed by Yamada--Watanabe rather than by the Lipschitz theory, and that strict positivity of $v$ 
requires the Feller condition $2\kappa\theta \geq \xi^2$ of Proposition \ref{prop:feller}, a 
condition frequently violated by calibrated parameters. Unlike local volatility, stochastic volatility renders the market \textbf{incomplete}. The filtration is generated by 
two Brownian motions while only one asset is traded, so the hypothesis of Theorem \ref{thm:mrt} 
fails: a claim whose value depends on $\Wt^v$ cannot be written as a stochastic integral against the 
traded asset alone. By Theorem \ref{thm:secondFTAP} the equivalent martingale measure is therefore 
not unique, and pricing requires either specifying a market price of volatility risk or, as is usual in practice, calibrating 
$(\kappa,\theta,\xi,\rho,v_0)$ directly to liquid quotes.

% The reason Heston's model, rather than some other stochastic volatility model, became the industry 
% standard is computational, and worth stating explicitly because Part \ref{part:part2} will build 
% directly on it. The characteristic function of the log-price is available in closed form.

\begin{theorem}[Heston's characteristic function]\label{thm:hestoncharfn}
Under \eqref{eq:heston}, write $\tau = T-t$ and $\Fwd_s = \Spot_s e^{(\rf-\divy)(T-s)}$ for the forward to $T$. Then
\begin{equation}\label{eq:hestoncf}
\charf(u,\tau) := \Ep{\QQ}\Bigl[\exp\Bigl(iu\log\tfrac{\Fwd_T}{\Fwd_t}\Bigr)\,\Big|\,\Ft\Bigr] 
= \exp\Bigl( C(u,\tau)\,\theta + D(u,\tau)\,v_t \Bigr),
\end{equation}
the characteristic function of the \emph{centred log-forward} $\log(\Fwd_T/\Fwd_t)$, where $C$ and $D$ solve the system of Riccati ordinary differential equations
\begin{equation}\label{eq:hestonriccati}
\frac{\partial D}{\partial \tau} = \frac{\xi^2}{2}D^2 + (\rho\xi i u - \kappa)D - \frac{u^2+iu}{2}, 
\qquad
\frac{\partial C}{\partial \tau} = \kappa D,
\end{equation}
with $C(u,0)=D(u,0)=0$. The system \eqref{eq:hestonriccati} admits an explicit solution in terms of 
elementary functions. \hfill $\square$
\end{theorem}

% The centring in \eqref{eq:hestoncf} is deliberate and should not be skipped over. Taking the
% characteristic function of $\log(\Fwd_T/\Fwd_t)$ rather than of the absolute log-spot $X_T$ strips
% out the factor $e^{iu\log \Fwd_t}$, and does so for two reasons. The first is arithmetical: the
% inversion formula below carries $\Fwd_t$ in its prefactor and again inside the log-moneyness
% $\lmon = \log(\Strike/\Fwd_t)$, so a characteristic function still carrying $\Fwd_t$ would introduce
% it a third time and the normalisation would be applied once and undone once. The second is
% structural: the centred $\charf$ depends only on $(v_t,\kappa,\theta,\xi,\rho,\tau)$ and not on
% $\Spot_t$, $\rf$ or $\divy$, which is exactly the separation between model parameters and market data
% that makes calibration well posed. This is the same reason the forward was singled out as the right
% variable in Subsection \ref{subsec:dividends}.

% The centring also supplies a cheap correctness check, worth recording because the error it catches is
% silent. Since $\Fwd$ is a $\QT$-martingale, $\Ep{\QQ}[\Fwd_T/\Fwd_t \mid \Ft] = 1$, so
% \begin{equation}\label{eq:hestoncheck}
% \charf(-i,\tau) = 1 \qquad \text{for every } \tau \geq 0 ,
% \end{equation}
% an identity satisfied by the centred characteristic function and violated by the uncentred one. Any
% implementation of \eqref{eq:hestoncf} should be tested against \eqref{eq:hestoncheck} before it is
% tested against prices.

Given \eqref{eq:hestoncf}, European option prices follow by Fourier inversion: the call price is 
recovered from the characteristic function through an integral of Lewis--Carr--Madan type \cite{carrmadan1999},
\begin{equation}\label{eq:fourierinversion}
C(\Strike,T) = \Bond{t}{T}\left[ \Fwd_t - \frac{\sqrt{\Fwd_t \Strike}}{\pi}\int_0^{\infty} 
\mathrm{Re}\left( e^{iu\lmon}\,\charf\bigl(u - \tfrac{i}{2},\tau\bigr)\right)\frac{du}{u^2 + \tfrac14} \right],
\end{equation}
a one-dimensional integral evaluated by quadrature in microseconds, with $\lmon = \log(\Strike/\Fwd_t)$ 
the log-moneyness of Subsection \ref{subsec:localvol}. Calibration --- which requires
pricing hundreds of options at each step of an optimiser --- is thereby feasible, whereas a model
requiring Monte Carlo for each European quote would not be.

Since $\Fwd$ is a $\QT$-martingale, $\Ep{\QQ}[\Fwd_T/\Fwd_t \mid \Ft] = 1$, so the centred
characteristic function \eqref{eq:hestoncf} must satisfy
\begin{equation}\label{eq:hestoncheck}
\charf(-i,\tau) = 1 \qquad \text{for every } \tau \geq 0 .
\end{equation}
Any implementation of \eqref{eq:hestoncf} should be tested against \eqref{eq:hestoncheck}, and
\eqref{eq:fourierinversion} against Black-76 in the limit $\xi \to 0$, before either is tested
against prices: both errors are otherwise silent.

One implementation caveat deserves a line, since it is a genuine trap rather than a nuisance. The 
explicit solution of \eqref{eq:hestonriccati} involves a complex logarithm, and the naive branch 
choice makes $\charf$ discontinuous in $\tau$ once the integrand has wound sufficiently far, 
corrupting prices at long maturities while leaving short ones correct. The remedy is the standard 
one in the literature, namely to write the solution in the form whose principal branch is 
continuous; the phenomenon is documented in \cite{bergomi2016}.

% That the tractability of the model reduces to a system of Riccati equations is a point worth 
% holding onto. The affine structure of \eqref{eq:heston} --- drift and variance both affine functions 
% of the state --- is exactly what produces \eqref{eq:hestonriccati}, and it is a structural property 
% rather than a happy accident. When we come to rough volatility in Part \ref{part:part2}, the model 
% will retain this affine character but lose the Markov property, and \eqref{eq:hestonriccati} will be 
% replaced by a \emph{fractional} Riccati equation. The pricing machinery \eqref{eq:fourierinversion} 
% survives essentially unchanged; only the equation for $D$ becomes non-local in time. Recording the 
% classical case here is what will make that generalisation a short step rather than a fresh start.

\subsubsection*{What stochastic volatility does and does not fix}

Stochastic volatility repairs the principal defect of local volatility: because the volatility has 
its own noise, the smile moves with the spot instead of flattening, and forward-starting products 
are priced far more sensibly. The correlation $\rho$ generates skew and the vol-of-vol $\xi$ 
generates curvature, both dynamically rather than by construction.

It does not, however, fit the surface exactly --- five parameters cannot match hundreds of quotes --- 
and it fails in one specific and instructive place. Empirically, the implied volatility skew of index 
options behaves like $T^{\Hurst-1/2}$ for small $T$, with $\Hurst$ well below $\tfrac12$, becoming very steep at
short maturities. Any model driven by standard Brownian motion, Heston included, produces a 
short-maturity skew that flattens as $T\to0$, because over short horizons a diffusive variance has 
had no time to move. No choice of $(\kappa,\theta,\xi,\rho,v_0)$ repairs this: the failure is not one 
of calibration but of the driving noise. This is the empirical observation from which Part 
\ref{part:part2} begins.

\subsection{Path-Dependence and American Options}\label{subsec:american}

We conclude this section by briefly discussing the challenges posed by path-dependent derivatives, focusing on the most liquid and prominent example: American options. 

Because an American option may be exercised at any time up to maturity, its value depends on the law of the underlying at \emph{every} intermediate date, not merely at $T$. A single implied volatility --- which by Subsection \ref{subsec:impliedvol} encodes one number extracted from one European quote --- therefore does not determine it: infinitely many dynamics reproduce that one quote while disagreeing about the intermediate marginals, and hence about the exercise boundary. This distinction is critically important because options on \textit{single names} (i.e., individual corporate stocks) are almost exclusively American-style. Investors in single equities are reluctant to relinquish the flexibility of early exercise, particularly to capture discrete dividends or navigate corporate actions. In contrast, broad market indices are generally considered more stable and are traded via extremely liquid European options.

Before developing any machinery it is worth recording how much of the early-exercise premium is
genuinely there, because for one of the two basic contracts the answer is: none. The result is due
to Merton \cite{merton1973} and explains why dividends have kept reappearing in the
discussion above. Its statement refers forward to the optimal stopping problem
\eqref{eq:optimalstopping}, written out two subsections below; the reader willing to grant that an
American option is worth the supremum, over exercise times, of the expected discounted payoff will
lose nothing by reading it now.

\begin{proposition}[No early exercise for the American call]\label{prop:mertonnoearly}
Assume $\rf \geq 0$ and that the underlying pays no dividends before $T$. Then an American call is
never optimally exercised before maturity, and its value coincides with that of the European call
with the same strike and maturity.
\end{proposition}

\textbf{\textit{Proof:}} Write $C(t)$ for the European call value at time $t$. Under $\QQ$ the
discounted spot is a martingale, so for any $t \leq T$, using the conditional Jensen inequality
(Proposition \ref{prop:condjensenl2}, item i)) applied to the convex function $x \mapsto (x-\Strike)^+$,
\[
C(t) = e^{-\rf(T-t)}\,\Ep{\QQ}\bigl[(\Spot_T-\Strike)^+ \mid \Ft\bigr]
\;\geq\; e^{-\rf(T-t)}\bigl(\Ep{\QQ}[\Spot_T\mid\Ft]-\Strike\bigr)^{+}
= \bigl(\Spot_t - \Strike e^{-\rf(T-t)}\bigr)^{+},
\]
the last step because $\Ep{\QQ}[\Spot_T\mid\Ft] = \Spot_t e^{\rf(T-t)}$ in the absence of dividends.
Since $\rf \geq 0$ gives $\Strike e^{-\rf(T-t)} \leq \Strike$, the right-hand side is at least
$(\Spot_t - \Strike)^{+}$, which is the immediate-exercise payoff. Holding therefore dominates
exercising at every $t < T$.

It remains to convert this pointwise domination into a statement about stopping times, which does
not follow from it directly. The chain above, read at a general $t$, says precisely that the
discounted payoff process $Z_t := e^{-\rf t}(\Spot_t-\Strike)^+$ satisfies
$\Ep{\QQ}[Z_T \mid \Ft] \geq Z_t$; that is, $Z$ is a $\QQ$-\emph{sub}martingale --- as it must be,
being a convex function of a martingale, by Proposition \ref{prop:condjensenl2}, item i), together
with $\rf \geq 0$. Doob's optional stopping theorem (Theorem \ref{thm:optionalstopping}, case i --- every stopping time being bounded by $T$ under Definition \ref{stoppingTime}) applied to a
submartingale gives, for every $\tau \in \Stop{[t,T]}$,
\[
\Ep{\QQ}\bigl[ Z_T \,\big|\, \mathcal{F}_\tau \bigr] \;\geq\; Z_\tau .
\]
Since $\tau \geq t$ we have $\Ft \subseteq \mathcal{F}_\tau$, so conditioning both sides on $\Ft$ and applying the tower property (Proposition \ref{prop:condexpproperties}, item 5) gives
\[
\Ep{\QQ}\bigl[e^{-\rf\tau}(\Spot_\tau-\Strike)^+ \,\big|\, \Ft\bigr]
\;\leq\; \Ep{\QQ}\bigl[e^{-\rf T}(\Spot_T-\Strike)^+ \,\big|\, \Ft\bigr] ,
\]
so the supremum in \eqref{eq:optimalstopping} is attained at $\tau = T$. \hfill $\blacksquare$

\medskip

Two things are worth extracting. The inequality is strict whenever $\rf > 0$ and the option has time
value, so the American call is not merely no better than the European one but strictly worse to
exercise early: one would be surrendering the interest earned on the strike. And the proof uses
$\Ep{\QQ}[\Spot_T\mid\Ft] = \Spot_t e^{\rf(T-t)}$, which is exactly what a dividend breaks. Once
$\divy > 0$, the forward is $\Spot_t e^{(\rf-\divy)(T-t)}$ and the chain of inequalities fails as
soon as $\divy > \rf$; early exercise then becomes genuinely optimal in a neighbourhood of a
dividend date. For the American \emph{put} the argument fails at the first step regardless of
dividends, since deep in the money the discounted payoff is bounded above while immediate exercise
is not, and early exercise is optimal for low enough spot (note that we do not have a Call-Put parity for American options).

We aim to show how a local volatility model can be adapted to price American options. However, we must first note a fundamental limitation: Local Volatility models do not perform perfectly for all types of path-dependent derivatives. In particular, because local volatility models make the volatility a deterministic function of the spot, hence perfectly correlated with it, they often misprice derivatives with discontinuous payoffs (which often lead to ill-posed pricing problems) or derivatives whose strikes are set at a future date, such as forward-starting options and cliquets. 

Nevertheless, to price American options accurately we must model the entire volatility surface prior 
to maturity, not merely the terminal marginal. The model we use, expressed under the $T$-forward 
measure of Subsection \ref{subsec:dividends}, is
\begin{equation}\label{eq:localvolforward}
    d\Fwd_t = \lvol(t,\Fwd_t)\,\Fwd_t \, d\Wt^{T}_t .
\end{equation}
% Working in the forward has the same three advantages set out in Subsection \ref{subsec:dividends}: 
% the drift vanishes because a forward is a martingale under its own measure, the dynamics stay 
% purely diffusive, and --- crucially here --- the discontinuities introduced by discrete cash 
% dividends are absorbed into the cash coefficient $b_t$ of \eqref{eq:spotforwardfactor} 
% and never enter \eqref{eq:localvolforward} at all.

\subsubsection*{Optimal stopping and the Snell envelope}

Before writing any partial differential equation, let us state what an American option \emph{is} in 
probabilistic terms, since the PDE formulation is derived from it and is opaque without it.

The holder may exercise at any stopping time of her choosing, and will choose the one maximising the 
value of the resulting cash flow. Writing $\Stop{[t,T]}$ for the set of stopping times valued in 
$[t,T]$ and $\PayoffFn$ for the exercise payoff, the value at time $t$ is
\begin{equation}\label{eq:optimalstopping}
V_t = \operatorname*{ess\,sup}_{\tau \in \Stop{[t,T]}} 
\Ep{\QQ}\left[ e^{-\int_t^{\tau}\rf(s)\,ds}\,\PayoffFn(\tau, \Spot_\tau) \,\Big|\, \Ft \right].
\end{equation}
The essential supremum, defined in Appendix \ref{app:esssup}, is required because the family is uncountable; this is an optimal stopping 
problem, and its solution is classical.

Discounting has to be brought inside before the classical statement applies, since it is the
\emph{discounted} value, not $V$ itself, that is a supermartingale. Write
\begin{equation}\label{eq:snellnormalisation}
Y_t := e^{-\int_0^{t}\rf(s)\,ds}\,\PayoffFn(t,\Spot_t), \qquad
\Snell_t := e^{-\int_0^{t}\rf(s)\,ds}\,V_t ,
\end{equation}
so that \eqref{eq:optimalstopping} becomes
$\Snell_t = \operatorname*{ess\,sup}_{\tau\in\Stop{[t,T]}} \Ep{\QQ}[\,Y_\tau \mid \Ft\,]$.

\begin{theorem}[Snell envelope]\label{thm:snell}
Let $Y$ be a right-continuous adapted process of class (D) --- Definition \ref{def:classD}. Then
the process $\Snell$ of \eqref{eq:snellnormalisation} is the \textbf{Snell envelope} of
$Y$: that is, $\Snell$ is the smallest càdlàg supermartingale dominating $Y$. It satisfies:
\begin{itemize}
    \item[i)] $\Snell_t \geq Y_t$ for all $t$ (domination);
    \item[ii)] $\Snell$ is a $\QQ$-supermartingale;
    \item[iii)] $\Snell$ is a $\QQ$-martingale on the \emph{continuation region}
    $\mathfrak{C} = \{(t,\omega) : \Snell_t > Y_t\}$;
\end{itemize}
and the stopping time $\tau^{*} = \inf\{ s \geq t : \Snell_s = Y_s \}$ is optimal in
\eqref{eq:optimalstopping}. A full treatment is given in \cite{peskirshiryaev2006}. \hfill $\square$
\end{theorem}

Since $e^{-\int_0^t \rf(s)\,ds}$ is strictly positive and deterministic, the continuation region and
the optimal stopping rule are the same whether read off $\Snell$ or off $V$; the normalisation
\eqref{eq:snellnormalisation} is bookkeeping, not a change of problem.

Property i) says the option is never worth less 
than exercising it. Property ii) says that holding the option is, on average and after discounting, a losing 
proposition: waiting has a cost, which is what makes early exercise ever attractive. Property iii) 
says that this cost is zero exactly where continuation is optimal, so that in the region where one 
does not exercise, the option is a fair game. The optimal rule is then transparent: continue while 
$\Snell > Y$, and stop the instant they meet.

\subsubsection*{The linear complementarity problem (LCP)}

Translating Theorem \ref{thm:snell} into the language of partial differential equations gives a free 
boundary problem, most rigorously expressed as a linear complementarity problem. Write
$\mathfrak{L}$ --- and not $\mathcal{L}$, which denotes the class of integrands of Subsection
\ref{subsec:itointegral} --- for the pricing operator associated with \eqref{eq:localvolforward},
\[
\mathfrak{L}V := \frac{\partial V}{\partial t} + \frac{1}{2}\lvol^2(t,\Fwd)\,\Fwd^2 \frac{\partial^2 V}{\partial \Fwd^2} - \rf(t)\,V .
\]

Exercise delivers the \emph{spot}, not the forward: the holder of an American put
receives $\Strike - \Spot_\tau$, and by \eqref{eq:spotforwardfactor} we have
$\Spot_t = a_t^{-1}(\Fwd_t - b_t)$. Expressed in the forward variable, the obstacle is therefore
\begin{equation}\label{eq:americanobstacle}
\PayoffFn(t,\Fwd) = \Bigl(\Strike - \tfrac{\Fwd - b_t}{a_t}\Bigr)^{+} \ \text{ for a put}, \qquad
\PayoffFn(t,\Fwd) = \Bigl(\tfrac{\Fwd - b_t}{a_t} - \Strike\Bigr)^{+} \ \text{ for a call},
\end{equation}
which is explicitly time-dependent through $a_t$ and $b_t$.

% Writing $(\Strike-\Fwd)^+$ instead, 
% as is sometimes done, silently assumes $a_t \equiv 1$ and $b_t \equiv 0$, and is simply wrong: it misplaces the
% exercise boundary, and it does so by the widest margin in a name paying large discrete dividends --- precisely 
% the case where early exercise matters most. Keeping $a_t$ and $b_t$ explicit costs nothing and preserves 
% validity for time-dependent $\rf(t)$ and $\divy(t)$, for cash and proportional dividends, and for 
% any mixture of them.

\begin{theorem}[American pricing as an LCP]\label{thm:LCP}
The value $V(t,\Fwd)$ of the American option satisfies, almost everywhere on 
$[0,T)\times(0,\infty)$, the three simultaneous conditions
\begin{equation}\label{eq:LCP_American}
    \begin{cases}
        \mathfrak{L}V \le 0 & \text{(supermartingale property, Theorem \ref{thm:snell}(ii))} \\[6pt]
        V(t, \Fwd) \ge \PayoffFn(t,\Fwd) & \text{(domination, Theorem \ref{thm:snell}(i))} \\[6pt]
        \bigl(\mathfrak{L}V\bigr) \cdot \bigl(V(t, \Fwd) - \PayoffFn(t,\Fwd)\bigr) = 0 & \text{(complementarity, Theorem \ref{thm:snell}(iii))}
    \end{cases}
\end{equation}
together with the terminal condition $V(T,\Fwd) = \PayoffFn(T,\Fwd)$ and $\PayoffFn$ as in 
\eqref{eq:americanobstacle}. \hfill $\square$
\end{theorem}

The correspondence with Theorem \ref{thm:snell} is exact. The complementarity condition encodes the dichotomy: 
at each point either $V > \PayoffFn$, in which case $\mathfrak{L}V=0$ and the ordinary pricing 
equation holds, or $V = \PayoffFn$ and the option is exercised. The boundary separating the two 
regions is not known in advance --- it is part of the solution --- which is what makes the problem 
free-boundary rather than merely parabolic, and what rules out any closed-form solution.

% \subsubsection{Calibration to Single Names}

% A practical hurdle arises at exactly this point. Dupire's formula (Theorem \ref{thm:dupire}) requires derivatives of \emph{European} prices, whereas single-name quotes are American, so the surface cannot be calibrated directly from the market. The standard remedy --- a short iterative mapping known as \textbf{de-Americanisation} --- is a calibration procedure rather than a piece of the theory, and we defer it to Section \ref{sec:deamericanisation} of Chapter \ref{ch:applications}, where the numerical machinery it needs is available.
\noindent\textbf{Sections \ref{sec:derivatives}--\ref{sec:othermodels}.}
The financial framework and pricing methodologies developed in the latter half of this chapter are built upon cornerstone texts in quantitative finance and stochastic analysis. The rigorous transition from abstract probability to continuous-time asset pricing, including the formulation of equivalent martingale measures and the \textit{No Free Lunch with Vanishing Risk} (NFLVR) condition, relies heavily on the definitive work by Delbaen and Schachermayer \cite{delbaenschachermayer1994}. This abstract measure-theoretic approach is bridged with continuous-time applications in the classic textbooks by Shreve \cite{shreve2004}, Øksendal \cite{oksendal2003}, and Björk \cite{bjork2009}.

The explicit formulation of diffusion models is rooted in the foundational breakthroughs of Black and Scholes \cite{blackscholes1973} and Merton \cite{merton1973}. To address the empirical shortcomings of this framework, Local Volatility was formalized by Dupire \cite{dupire1994}, with its practical translation into the implied volatility domain extensively covered by Gatheral \cite{gatheral2006}. For an exhaustive, practitioner-oriented treatment of volatility smile dynamics and Stochastic Volatility models (such as the Heston model \cite{heston1993}), Bergomi \cite{bergomi2016} serves as an indispensable reference. Finally, the Fourier pricing route of \eqref{eq:fourierinversion}, which is what makes calibrating an affine model such as Heston computationally feasible at all, originates with Carr and Madan \cite{carrmadan1999}; the particular form used here, symmetric in $\Fwd$ and $\Strike$ and integrated along the line $\mathrm{Im}\,u = -\tfrac12$, is the variant that is numerically best behaved, and the arbitrage-free interpolation it is calibrated against is Gatheral's SVI \cite{gatheral2014}, with the wing bounds supplied by Lee's moment formula \cite{lee2004}.
